\chapter{Folgen, Grenzwerte und Cauchy-Eigenschaft}
\label{b47guide:ch:sequences}

\section{Metrische Konvergenz}

\begin{definition}[Konvergenz und Teilfolge]
Eine Folge $(x_n)_{n\geq0}$ in $X$ \emph{konvergiert} gegen $x\in X$, wenn
\[
 \forall\varepsilon>0\ \exists N\geq0\ \forall n\geq N:
 d(x_n,x)<\varepsilon.
\]
Wir schreiben $x_n\to x$. Eine \emph{Teilfolge} hat die Gestalt $(x_{n_k})$
mit streng wachsenden natürlichen Indizes $n_k$.
\end{definition}

\begin{theorem}[Grenzwertgesetze ohne algebraische Struktur]
Ein metrischer Grenzwert ist eindeutig. Jede gegen $x$ konvergente Folge
ist beschränkt, und jede ihrer Teilfolgen konvergiert gegen $x$.
Endlich viele Änderungen und das Weglassen endlich vieler Anfangsglieder
ändern die Konvergenz und ihren Grenzwert nicht. Schließlich konstante
Folgen konvergieren gegen ihren schließlich angenommenen Wert.
\end{theorem}

\begin{theorem}[Grenzwerte von Abständen]
Aus $x_n\to x$ und $y_n\to y$ folgt $d(x_n,y_n)\to d(x,y)$.
Ist $d(x_n,y_n)\to0$ und $x_n\to x$, so gilt auch $y_n\to x$.
Eine Folge in $A\subseteq X$ konvergiert im Teilraum $A$ genau dann gegen
$a\in A$, wenn sie im Raum $X$ gegen $a$ konvergiert.
\end{theorem}

\begin{theorem}[Folgenkriterien für Abschluss und Abgeschlossenheit]
\label{b47guide:thm:sequential-closure}
Für $A\subseteq X$ und $x\in X$ gilt
\[
 x\in\overline A
 \quad\Longleftrightarrow\quad
 \text{es gibt eine Folge $(a_n)$ in $A$ mit $a_n\to x$.}
\]
Damit ist $A$ genau dann abgeschlossen, wenn jeder in $X$ existierende
Grenzwert einer Folge aus $A$ wieder in $A$ liegt.
\end{theorem}

\begin{definition}[Teilfolgengrenzwert und Häufungspunkt einer Menge]
$x\in X$ heißt \emph{Teilfolgengrenzwert} von $(x_n)$, wenn eine Teilfolge
gegen $x$ konvergiert. $x$ heißt \emph{Häufungspunkt} einer Menge
$A\subseteq X$, wenn jede Kugel um $x$ einen Punkt von $A\setminus\{x\}$
enthält. Diese beiden Begriffe unterscheiden sich bei wiederholten
Folgenwerten.
\end{definition}

\begin{theorem}[Kriterien für Häufung]
$x$ ist genau dann Teilfolgengrenzwert von $(x_n)$, wenn jede Kugel um $x$
unendlich viele Folgenglieder, nach Indizes gezählt, enthält.
$x$ ist genau dann Häufungspunkt von $A$, wenn eine Folge paarweise
verschiedener Punkte aus $A\setminus\{x\}$ gegen $x$ konvergiert.
\end{theorem}

\section{Cauchy-Folgen}

\begin{definition}[Cauchy-Folge]
\label{b47guide:def:cauchy}
Eine Folge $(x_n)$ heißt \emph{Cauchy-Folge}, wenn
\[
 \forall\varepsilon>0\ \exists N\geq0\ \forall m,n\geq N:
 d(x_m,x_n)<\varepsilon.
\]
Die Bedingung verlangt keinen bereits benannten Grenzwert.
\end{definition}

\begin{theorem}[Grundlegende Cauchy-Kriterien]
Jede konvergente Folge ist eine Cauchy-Folge. Jede Cauchy-Folge ist
beschränkt, und jede ihrer Teilfolgen ist eine Cauchy-Folge. Endliche
Änderungen und das Weglassen endlich vieler Anfangsglieder erhalten die
Cauchy-Eigenschaft. Für Folgen in einem Teilraum stimmt die Cauchy-Bedingung
im Teilraum mit derjenigen im umgebenden Raum überein.
\end{theorem}

\begin{theorem}[Cauchy-Folge mit einem Teilfolgengrenzwert]
\label{b47guide:thm:cauchy-subsequence}
Besitzt eine Cauchy-Folge eine gegen $x\in X$ konvergente Teilfolge,
so konvergiert bereits die ganze Folge gegen $x$.
\end{theorem}

\begin{theorem}[Asymptotisch gleiche Folgen]
Gilt $d(x_n,y_n)\to0$, so ist $(x_n)$ genau dann eine Cauchy-Folge, wenn
$(y_n)$ eine Cauchy-Folge ist. Sind $(x_n)$ und $(y_n)$ beide Cauchy-Folgen,
so besitzt die reelle Folge $d(x_n,y_n)$ einen endlichen Grenzwert.
\end{theorem}

\begin{theorem}[Summierbare Schrittweiten]
Gilt $d(x_n,x_{n+1})\leq a_n$ mit $a_n\geq0$ und
$\sum_{n=0}^{\infty}a_n<\infty$, so ist $(x_n)$ eine Cauchy-Folge.
Ist $X$ vollständig, so existiert ein Grenzwert $x\in X$ und
\[
 d(x_n,x)\leq\sum_{k=n}^{\infty}a_k.
\]
\end{theorem}

\begin{theorem}[Geometrisch kleine Schritte als Vollständigkeitskriterium]
$X$ ist genau dann vollständig, wenn jede Folge mit
$d(x_n,x_{n+1})\leq2^{-n}$ in $X$ konvergiert.
\end{theorem}